\documentclass[12pt]{article}
\usepackage[margin=1.15in]{geometry}
\usepackage{amsmath,amssymb,amsthm}
\usepackage{graphicx}
\usepackage[hyphens]{url}
\usepackage[colorlinks=true,citecolor=blue,linkcolor=blue,urlcolor=blue,breaklinks=true]{hyperref}

\newtheorem{theorem}{Theorem}
\newtheorem{proposition}{Proposition}
\newtheorem{corollary}{Corollary}
\theoremstyle{definition}
\newtheorem{assumption}{Assumption}
\theoremstyle{remark}
\newtheorem{remark}{Remark}

\newcommand{\E}{\mathbb{E}}
\newcommand{\bth}{\bar\theta}
\newcommand{\D}{\mathbb{D}}

\title{The Price of Risk When Aggregation Fails:\\
Relative Wealth Concerns with Heterogeneous Risk Aversion}
\author{Renato Vicente\thanks{Dept.\ of Applied Mathematics, Instituto de
Matem\'atica, Estat\'istica e Ci\^encia da Computa\c{c}\~ao, Universidade de
S\~ao Paulo. \texttt{rvicente@ime.usp.br}}}
\date{\today}

\begin{document}
\maketitle

\begin{abstract}
Gómez (2007) shows that keeping-up-with-the-Joneses preferences with linear risk
tolerance collapse into a rescaling of the representative agent's risk aversion,
but only under restrictive conditions requiring common preference parameters.
This paper characterises the failure, in a CARA economy where each investor
compares his wealth to a weighted average over a network of peers. With common
risk aversion the comparison is exactly neutral, for every influence network and
every profile of comparison intensities, extending Gómez's CARA case to
networks. With heterogeneous risk aversion the demand scales are a
Katz--Bonacich centrality and the effective risk aversion depends on the
network. In a mean-field limit it is $H(\theta+\lambda)-\lambda$,
$H$ the harmonic mean, so comparison interpolates aggregation from the harmonic
to the arithmetic mean at a rate equal to the squared coefficient of variation
of shifted risk tolerance. Away from that limit it need not be monotone: in a
three-investor network the coefficient first rises, then falls below its
no-comparison value. Reversal becomes rare as the population grows, but this
reflects mixing rather than size. And homophily in risk aversion can suppress
the effect entirely, the harmonic mean surviving at arbitrarily large comparison
intensity.
\end{abstract}

\noindent\textbf{Keywords:} relative wealth concerns, keeping up with the
Joneses, CARA aggregation, Katz--Bonacich centrality, homophily.

\medskip
\noindent\textbf{JEL:} G11, G12, D53, D85.

%=====================================================================
\section{Introduction}

\cite{gomez} studies an exchange economy in which agents evaluate consumption
relative to a weighted average of the consumption of others, and asks when
equilibrium prices depend only on aggregates. His Proposition~1 shows that
linear risk tolerance with common cautiousness delivers a representative agent
and rules out portfolio bias; his Proposition~2 gives the conditions under
which prices depend on the aggregate endowment alone, and finds them
restrictive: the preference parameter must be common across agents and, absent
the trivial case of uniform weights, so must the equilibrium value of the
endowment. Under those conditions the CARA case reduces cleanly, and keeping up
with the Joneses is equivalent to lowering the representative agent's risk
aversion from $\alpha$ to $\alpha(1-\gamma)$.

What happens outside those conditions is left open, and it is the natural
question. \cite{gomez} notes that in their absence equilibrium prices may be
affected by the distribution of weights in the definition of the Joneses, by
the preference parameters, and by the endowment distribution --- three channels
that a representative-agent reduction cannot express. This paper opens them,
in the CARA case, for a benchmark that is agent-specific: investor $a$ compares
himself to $\sum_jA_{aj}W_j$, a weighted average over a network of influence
rather than a single economy-wide aggregate.

The results divide according to whether risk aversion is common.

When it is, the comparison is exactly neutral. Every investor holds the
position he would hold with no relative concern at all, for every
row-stochastic $A$ and every profile of comparison intensities $\lambda_a$
(Theorem~\ref{thm:neutral}). This is the network extension of Gómez's CARA
case, and its content is that the network drops out entirely: no statistic of
$A$ survives into the equilibrium. It is worth stating separately because the
proof is immediate and because it delimits precisely where the interesting
region begins.

When risk aversion is heterogeneous, neutrality fails and the network matters.
The demand scales solve $[\D_{\theta+\lambda}-\D_{\lambda}A]\vec c=\mathbf 1$,
a Katz--Bonacich system in the sense of \cite{ballester}, and we study it in
three ways. Under a mean-field condition the effective market risk aversion is
$H(\theta+\lambda)-\lambda$, so that comparison interpolates the aggregation of
risk aversion from the harmonic mean at $\lambda=0$ to the arithmetic mean as
$\lambda\to\infty$, at a rate equal to the squared coefficient of variation of
shifted risk tolerance (Proposition~\ref{prop:mf}). Away from that condition
the direction is not guaranteed and the relation need not be monotone: a
three-investor network with a directed pattern of attention delivers a
coefficient that first rises and then falls below its no-comparison value
(Proposition~\ref{prop:reversal}). We measure how often that occurs under a
dense i.i.d.-weight design, where it becomes rare as the population grows; this
is a statement about mixing rather than about size, and block structure
preserves reversal at any number of investors. And if investors compare themselves only to others of their
own type, the aggregation does not move at all, however intense the comparison
(Proposition~\ref{prop:homophily}); homophily in risk aversion blocks the
mechanism completely.

Together these say that the comparison channel in this class of models is not a
property of preferences but of the interaction between preference dispersion
and network structure, and that the mean-field intuition --- which is what a
representative-agent treatment would deliver --- can fail in both magnitude and
sign.

%=====================================================================
\section{Related literature}

Preferences over relative consumption or wealth have been used to address a
range of asset-pricing questions. \cite{abel} and \cite{gali} build
consumption-based models with external benchmarks and obtain an amplified
equity premium; \cite{campbell} show how an external habit generates
time-varying effective risk aversion; \cite{chankogan} make the wealth
distribution a state variable under heterogeneous CRRA preferences, which is
the dynamic counterpart of the configuration studied here. \cite{dupor}
provide the taxonomy under which \eqref{eq:pref} is jealousy combined with
keeping up with the Joneses.

Within the CARA class, \cite{degennaro} study a two-agent economy with multiple
risky assets and different risk tolerances, showing that relative concerns are
equivalent to modified tolerances and lower equilibrium premia; that economy is
Section~\ref{sec:hetero} at $m=2$. \cite{demarzo04} obtain portfolio bias when
the benchmark differs in composition from the investor's own holdings, a
channel excluded here by the scalar structure of \eqref{eq:foc}.
\cite{lacker} solve relative-performance games for CARA and CRRA investors in
mean-field and $n$-agent form; their concern is the existence and structure of
Nash equilibria rather than whether the aggregate moves, and their mean-field
formulation is the analogue of our Proposition~\ref{prop:mf} in a different
setting. \cite{kraft} take the problem to incomplete markets in continuous
time. \cite{panageas} surveys what heterogeneity and inequality imply for asset
prices more broadly.

Two connections deserve care. \cite{gollier} gives the criterion governing the
effect of endowment inequality on the equity premium, namely concavity of risk
tolerance in wealth. Our \eqref{eq:cv2} is an analogue for dispersion in a
preference parameter rather than in endowments, and should be read as such
rather than as an application of his theorem. And \cite{ballester} give the
linear-quadratic network game whose solution \eqref{eq:system} reproduces; what
the present setting adds is that the centrality vector becomes constant, and
every statistic of the network drops out, precisely when risk aversion is
common.

%=====================================================================
\section{The economy}

There are $m$ investors, one riskless asset with net return $r_f$, and $n$
risky assets with excess returns $\vec y$ that are Gaussian with mean
$\vec r_e$ and non-singular covariance $\Sigma$. Investor $a$ begins with
wealth $W_a^0$ and chooses monetary positions $\vec\pi_a\in\mathbb{R}^n$, so
that $W_a=W_a^0(1+r_f)+\vec\pi_a^{\top}\vec y$. Preferences are
%
\begin{equation}
U_a=-\exp\Big[-\theta_a W_a-\lambda_a\Big(W_a-\sum_jA_{aj}W_j\Big)\Big],
\qquad \theta_a>0,\ \lambda_a\geq0 .
\label{eq:pref}
\end{equation}

\begin{assumption}[Normalised benchmark]
\label{ass:norm}
$A_{aa}=0$, $A_{aj}\geq0$ and $\sum_jA_{aj}=1$ for every $a$.
\end{assumption}

The benchmark is thus a weighted average of the wealth of the other investors,
as in \cite{gomez}, with the difference that the weights may vary across
investors. Assumption~\ref{ass:norm} makes the preference invariant to a
uniform translation of all wealths: adding a constant to every $W_j$ shifts the
exponent by $-\theta_a$ times that constant, independently of $\lambda_a$.
It is also what the results below require. Dropping it changes the object
being compared: row sums below one describe comparison with a partial or
external benchmark, which is a coherent alternative we do not study, while row
sums above $1+\theta_a/\lambda_a$ make an investor worse off when the whole
economy becomes richer, which is harder to motivate. We retain
Assumption~\ref{ass:norm} throughout and claim nothing about either case.

\begin{assumption}[Nash behaviour]
\label{ass:nash}
Each investor chooses $\vec\pi_a$ taking the positions $\vec\pi_j$, $j\neq a$,
as fixed at their equilibrium values. An equilibrium is a profile
$(\vec\pi_1,\dots,\vec\pi_m)$ and a vector $\vec r_e$ such that each
$\vec\pi_a$ is optimal given the others and markets clear.
\end{assumption}

Assumption~\ref{ass:nash} is the standard notion and we state it explicitly
because the conjecture it embodies matters for what follows. An investor
behaving in this way conjectures that the positions of others do not respond to
his own: the conjectural derivative is zero, and the conjectured levels are the
equilibrium positions. This is distinct from an investor who believes that
others hold whatever he holds, a possibility we take up separately in
Section~\ref{sec:projection}.

\subsection*{First-order conditions}

Under \eqref{eq:pref} the random part of the exponent is
$-\big[(\theta_a+\lambda_a)\vec\pi_a-\lambda_a\sum_jA_{aj}\vec\pi_j\big]^{\top}\vec y$.
Writing $\vec q_a$ for the bracket, Gaussian returns give
$\E[U_a]=-C\exp\big(-\vec q_a^{\top}\vec r_e+\tfrac12\vec q_a^{\top}\Sigma\vec q_a\big)$,
whose exponent is strictly convex in $\vec\pi_a$ because
$\partial\vec q_a/\partial\vec\pi_a=(\theta_a+\lambda_a)\mathbb{I}$ is
non-singular. The optimum is interior and solves $\vec q_a=\Sigma^{-1}\vec r_e$,
that is
%
\begin{equation}
(\theta_a+\lambda_a)\,\vec\pi_a-\lambda_a\sum_jA_{aj}\vec\pi_j
=\Sigma^{-1}\vec r_e .
\label{eq:foc}
\end{equation}
%
Every solution is proportional to $\Sigma^{-1}\vec r_e$, so all investors hold
the same portfolio of risky assets and differ only in scale; two-fund
separation holds, as in Gómez's Proposition~1. Writing
$\vec\pi_a=c_a\Sigma^{-1}\vec r_e$ reduces \eqref{eq:foc} to
%
\begin{equation}
\big[\D_{\theta+\lambda}-\D_{\lambda}A\big]\,\vec c=\mathbf 1,
\qquad \D_x=\mathrm{diag}(x_1,\dots,x_m),
\label{eq:system}
\end{equation}
%
whose solution is a Katz--Bonacich centrality of the influence network
\cite{ballester}. The matrix is strictly diagonally dominant for every
positive $\theta_a$, since the off-diagonal row sums are
$\lambda_a\sum_jA_{aj}=\lambda_a<\theta_a+\lambda_a$; existence and uniqueness
therefore hold, and $\vec c>0$ follows from the Neumann expansion. Market
clearing against a supply with value weights $\vec w$ and aggregate wealth
$W_M^0$ gives
%
\begin{equation}
\vec r_e=\theta_MW_M^0\Sigma\vec w,
\qquad \frac{1}{\theta_M}=\sum_ac_a,
\qquad \bth=m\,\theta_M .
\label{eq:clear}
\end{equation}
%
We call $\bth$ the effective market risk-aversion coefficient and reserve
\emph{mean--covariance price-of-risk coefficient} for the product
$W_M^0\theta_M=\bar W^0\bth$, which is the dimensionless object; $\bth$ alone
carries units of inverse wealth and we do not call it a price.

\begin{remark}[What the model determines]
Equations \eqref{eq:foc}--\eqref{eq:clear} determine a mean--covariance
relation given $\Sigma$ and the supplies. They are not a complete system of
asset prices: with payoffs and quantities as primitives a change in prices
would move both the mean and the covariance of returns. Statements below
concern demands and expected excess returns conditional on $\Sigma$, and
nothing is claimed about price levels, market volatility or the invariance of
betas.
\end{remark}

%=====================================================================
\section{Neutrality under common risk aversion}

\begin{theorem}
\label{thm:neutral}
Let Assumptions~\ref{ass:norm} and \ref{ass:nash} hold and suppose
$\theta_a=\theta$ for every $a$. Then $\vec c=\theta^{-1}\mathbf 1$ is the
unique solution of \eqref{eq:system}, so every investor holds
$\vec\pi_a=\Sigma^{-1}\vec r_e/\theta$ and $\bth=\theta$, for every influence
matrix satisfying Assumption~\ref{ass:norm} and every profile
$(\lambda_1,\dots,\lambda_m)$.
\end{theorem}

\begin{proof}
Substituting $\vec c=\theta^{-1}\mathbf 1$ in row $a$ of \eqref{eq:system} and
using $\sum_jA_{aj}=1$ gives
$(\theta+\lambda_a)/\theta-\lambda_a/\theta=1$. Uniqueness follows from strict
diagonal dominance, as above.
\end{proof}

The cancellation occurs row by row, before any aggregation, which is why
nothing about $A$ or about the dispersion of $\lambda_a$ enters. Economically, the
benchmark is fixed in a unilateral deviation: under Assumption~\ref{ass:nash}
the positions of others do not respond. Neutrality arises instead because at
the equilibrium profile the benchmark carries the same risky exposure as the
investor himself, so that
$(\theta+\lambda_a)\vec\pi-\lambda_aA\vec\pi=(\theta+\lambda_a)\vec\pi-\lambda_a\vec\pi=\theta\vec\pi$.
The benchmark exposure offsets the additional $\lambda_a$ loading on own
equilibrium wealth in the level of total risky exposure, not through any
conjectured response to the deviation.

Theorem~\ref{thm:neutral} is the network extension of the CARA case of
\cite{gomez}, and we claim no more for it than that. His Proposition~2
identifies common cautiousness --- here, common $\theta$ --- as what makes the
weighting drop out; the addition is that the weighting may be agent-specific
and the comparison intensities may differ, and neither disturbs the conclusion.

\begin{remark}[Relation to the keeping-up parametrization]
\label{rem:param}
The additive form \eqref{eq:pref} and the multiplicative form used by
\cite{gomez} describe the same preference. Collecting terms in the exponent,
%
\begin{equation}
-\theta_aW_a-\lambda_a\Big(W_a-\sum_jA_{aj}W_j\Big)
=-\alpha_a\Big(W_a-\gamma_a\sum_jA_{aj}W_j\Big),
\label{eq:param}
\end{equation}
%
with
$\alpha_a=\theta_a+\lambda_a$, $\gamma_a=\lambda_a/(\theta_a+\lambda_a)$ and
hence $\theta_a=\alpha_a(1-\gamma_a)$: $\alpha_a$ is the local coefficient of
absolute risk aversion in own wealth and $\gamma_a\in[0,1)$ the keeping-up
weight. In this notation \cite{gomez} writes the CARA case as
$u(c,X)=-\exp(-\alpha(c-\gamma X))$ and reduces it to a representative agent
with coefficient $\alpha(1-\gamma)$, which is exactly $\theta$; his conclusion
and Theorem~\ref{thm:neutral} therefore agree wherever both apply.

The parametrizations are not interchangeable as \emph{hypotheses}, however. His
Proposition~2 requires the preference parameters to be identical across agents
and, barring uniform weights, the equilibrium value of the endowment to be
identical as well. Theorem~\ref{thm:neutral} requires neither. It constrains
only the product $\alpha_a(1-\gamma_a)=\theta$, so neutrality holds on the whole
level set $\{(\alpha,\gamma):\alpha(1-\gamma)=\theta\}$, along which an investor
with a stronger comparison motive is also locally more risk averse in own
wealth; and initial wealth is absent from \eqref{eq:foc} altogether, since CARA
demands have no wealth effects, so no restriction on the endowment distribution
is needed. Reading $\theta_a$ rather than $\alpha_a$ as ``the'' risk
aversion is what makes the statement look stronger than it is, and it is
$\theta_a$ --- the coefficient net of the comparison --- whose dispersion drives
everything in Section~\ref{sec:hetero}.
\end{remark}

\begin{remark}[Two notions of neutrality]
\label{rem:two}
Theorem~\ref{thm:neutral} asserts that individual positions are unchanged,
which is stronger than saying the effective coefficient is unchanged. The latter
requires only $\sum_ac_a=m/\theta$, and deviations can offset. With four
investors, $\theta=1$ and demand scales
$\vec c=(\tfrac57,\tfrac97,1,1)$ --- corresponding to coefficients
$\kappa=(\tfrac75,\tfrac79,1,1)$ --- one has $\sum_ac_a=4$ and hence
$\bth=\theta$ exactly, although the first two investors hold positions
differing from the no-comparison benchmark by $-\tfrac27$ and $+\tfrac27$, or
about $28.6$ per cent in each direction. Results below concern $\bth$, so the weaker
notion is the relevant one, and statements about it should not be read as
statements about allocations.
\end{remark}

%=====================================================================
\section{Heterogeneous risk aversion}
\label{sec:hetero}

We now let $\theta_a$ vary, which is exactly the configuration in which the
aggregation conditions of \cite{gomez} fail. The system \eqref{eq:system} still
has a unique positive solution, but it is no longer proportional to
$\mathbf 1$, and $\bth$ depends on $\lambda$ and on $A$.

\subsection{The mean-field limit}

\begin{proposition}
\label{prop:mf}
Let $\lambda_a=\lambda$ for all $a$ and consider a sequence of economies
indexed by $m$ in which the empirical distribution of $\theta_a$ converges
weakly to a law with $\E[(\theta+\lambda)^{-1}]<\infty$ and
$\{(\theta+\lambda)^{-1}\}$ uniformly integrable, and in which
$\sum_jA_{aj}c_j-\bar c_m\to0$ uniformly in $a$, where
$\bar c_m=m^{-1}\sum_{a=1}^{m}c_a$. Then, writing $\bth_m$ for the effective
coefficient of the $m$-investor economy,
%
\begin{equation}
\bth_m(\lambda)\longrightarrow \bth_\infty(\lambda)=H(\theta+\lambda)-\lambda,
\qquad H(x)=\big(\E[1/x]\big)^{-1} .
\label{eq:mf}
\end{equation}
%
The limit $\bth_\infty$ is differentiable wherever the moments below exist,
and
%
\begin{equation}
\frac{d\bth_\infty}{d\lambda}
=\frac{\mathrm{Var}\big[(\theta+\lambda)^{-1}\big]}
{\E\big[(\theta+\lambda)^{-1}\big]^{2}}\;\geq\;0 ,
\label{eq:cv2}
\end{equation}
%
with equality if and only if $\theta$ is degenerate. If moreover
$\E[\theta^3]<\infty$ then
$\bth_\infty=\E[\theta]-\mathrm{Var}(\theta)/\lambda+O(\lambda^{-2})$ as
$\lambda\to\infty$. Equation \eqref{eq:cv2} is the derivative of the limit;
we do not claim it is the limit of the derivatives of the finite economies.
\end{proposition}

\begin{proof}
Under the mixing condition row $a$ of \eqref{eq:system} becomes
$(\theta_a+\lambda)c_a=1+\lambda\bar c$ in the limit, so
$c_a=(1+\lambda\bar c)/(\theta_a+\lambda)$. Averaging and writing
$M=\E[(\theta+\lambda)^{-1}]$ gives $\bar c=M/(1-\lambda M)$ and
$\bth=1/M-\lambda$. Differentiating gives \eqref{eq:cv2}, whose sign is
Cauchy--Schwarz, with equality only for degenerate $\theta$. The expansion
follows from
$\E[(\theta+\lambda)^{-1}]
=\lambda^{-1}\big(1-\E[\theta]/\lambda+\E[\theta^{2}]/\lambda^{2}
+O(\lambda^{-3})\big)$, which requires the third moment to control the
remainder.
\end{proof}

In this limit comparison interpolates the aggregation of risk aversion between
the two classical aggregators: the harmonic mean at $\lambda=0$, which is the
standard CARA result, and the arithmetic mean as $\lambda\to\infty$. Since
$H(\theta)\leq\E[\theta]$, the largest increase in the effective coefficient
available is the harmonic--arithmetic gap of the risk-aversion distribution, a
statistic of the cross-section rather than a free parameter.

\begin{remark}
The mixing condition is stated in terms of the endogenous vector $\vec c$, and
is therefore close to assuming the conclusion. A primitive sufficient condition
is that $A$ have rows that are identical up to the leave-one-out correction,
$A_{aj}=1/(m-1)$ for $j\neq a$; then $\sum_jA_{aj}c_j=(\sum_kc_k-c_a)/(m-1)$
differs from $\bar c$ by $O(1/m)$ and \eqref{eq:mf} holds in the limit. Exact
equality at finite $m$ would require identical rows, which
Assumption~\ref{ass:norm} forbids since $A_{aa}=0$. What conditions on $A$
short of this suffice is open, and in our view is the most useful direction in
which this analysis can be extended.
\end{remark}

\subsection{Non-monotonicity and sign reversal}

\begin{proposition}
\label{prop:reversal}
There exist normalised benchmarks and risk-aversion profiles for which
$d\bth/d\lambda<0$ on a non-empty interval and for which
$\bth(\lambda)<\bth(0)$ at sufficiently large comparison intensities.
\end{proposition}

\begin{proof}
Take $m=3$, $\theta=(\tfrac12,1,2)$, $\lambda_a=\lambda$ and
%
\begin{equation*}
A=\begin{pmatrix}0&1&0\\1&0&0\\1&0&0\end{pmatrix},
\end{equation*}
%
which satisfies Assumption~\ref{ass:norm}. Solving \eqref{eq:system} gives
%
\begin{equation}
\bth(\lambda)=\frac{3(\lambda+2)(3\lambda+1)}{12\lambda^{2}+24\lambda+7},
\qquad
\bth'(\lambda)=-\frac{3\big(12\lambda^{2}+6\lambda-1\big)}
{\big(12\lambda^{2}+24\lambda+7\big)^{2}} .
\label{eq:cex}
\end{equation}
%
Hence $\bth'(0)=3/49>0$ and $\bth'(\lambda)<0$ for
$\lambda>(\sqrt{21}-3)/12\approx0.132$. The function rises from
$\bth(0)=6/7\approx0.857$ to a maximum $\approx0.860$ at that point, then
falls, returning to $6/7$ at $\lambda=1/3$ and continuing downwards:
$\bth(1)=0.837209$ and $\bth(2)=0.815534$.
\end{proof}

The example is more informative than a monotone decrease would be. Comparison
first raises the effective coefficient, as the mean-field intuition of
Proposition~\ref{prop:mf} suggests, and only at higher intensity does the
structure of attention take over and reverse the direction. The pattern is what
produces it: the most risk averse investor looks at the least risk averse,
while the other two look only at each other.

Reversal is not generic. Under the dense i.i.d.-weight design described below
it becomes rare as the number of investors grows and local benchmark exposures
approach complete mixing; Remark~\ref{rem:blocks} shows that this is a property
of that design and not of size.
Table~\ref{tab:reversal} reports two distinct frequencies over $20{,}000$
independent draws at each size: the \emph{cumulative} event
$\bth(0.5)<\bth(0)$, and the \emph{marginal} event $\bth'(0.5)<0$. The
derivative is computed analytically rather than by finite differences: writing
$M(\lambda)=\D_{\theta}+\lambda(\mathbb{I}-A)$ and $\vec c=M^{-1}\mathbf 1$,
%
\begin{equation}
\vec c\,'=-M^{-1}(\mathbb{I}-A)\vec c,
\qquad
\bth'=-\frac{m\,\mathbf 1^{\top}\vec c\,'}{(\mathbf 1^{\top}\vec c)^{2}} .
\label{eq:deriv-num}
\end{equation}
%
The number of investors is held fixed within each row, so that the comparison
across rows identifies the effect of size. Off-diagonal entries of $A$ are
uniform on $(0,1)$ before row normalisation, and
$\theta_a=\max\{1+\sigma Z_a,\,0.15\}$ with $Z_a$ standard normal.

\begin{table}[htbp]
\centering
\setlength{\tabcolsep}{5pt}
\begin{tabular}{cc@{\hskip 12pt}cc@{\hskip 12pt}cc}
\hline
& & \multicolumn{2}{c@{\hskip 12pt}}{cumulative: $\bth(0.5)<\bth(0)$}
  & \multicolumn{2}{c}{marginal: $\bth'(0.5)<0$} \\
$m$ & $\sigma$ & frequency & $95\%$ interval & frequency & $95\%$ interval \\
\hline
3  & 0.4 & 0.0919 & $[0.0879,\,0.0959]$ & 0.0952 & $[0.0912,\,0.0993]$ \\
3  & 0.8 & 0.0340 & $[0.0316,\,0.0366]$ & 0.0372 & $[0.0346,\,0.0399]$ \\
6  & 0.4 & 0.0057 & $[0.0047,\,0.0068]$ & 0.0062 & $[0.0052,\,0.0073]$ \\
6  & 0.8 & 0.0008 & $[0.0005,\,0.0014]$ & 0.0008 & $[0.0005,\,0.0012]$ \\
10 & 0.4 & 0.0001 & $[0.0000,\,0.0004]$ & 0.0001 & $[0.0000,\,0.0004]$ \\
10 & 0.8 & 0.0000 & $[0.0000,\,0.0002]$ & 0.0000 & $[0.0000,\,0.0002]$ \\
20 & 0.4 & 0.0000 & $[0.0000,\,0.0002]$ & 0.0000 & $[0.0000,\,0.0002]$ \\
20 & 0.8 & 0.0000 & $[0.0000,\,0.0002]$ & 0.0000 & $[0.0000,\,0.0002]$ \\
60 & 0.4 & 0.0000 & $[0.0000,\,0.0002]$ & 0.0000 & $[0.0000,\,0.0002]$ \\
60 & 0.8 & 0.0000 & $[0.0000,\,0.0002]$ & 0.0000 & $[0.0000,\,0.0002]$ \\
\hline
\end{tabular}
\caption{Frequency of negative cumulative and marginal effects at
$\lambda=0.5$, over $20{,}000$ draws per row with Wilson intervals. All ten
cells of the design are reported; entries of $0.0000$ are exact zeros, no draw
out of $20{,}000$ producing the event, and $0.0002$ is the upper endpoint of the
$95\%$ Wilson interval when no event is observed. Under this design reversal is already below one
per cent at $m=6$ and indistinguishable from zero at $m\geq10$. It is also
rarer at higher dispersion of risk aversion, which is consistent with the
positive mean-field effect of \eqref{eq:cv2} growing with dispersion and so
becoming harder for network asymmetry to overturn. The decay in $m$ is specific
to the dense i.i.d.-weight ensemble; see Remark~\ref{rem:blocks}.}
\label{tab:reversal}
\end{table}

\begin{remark}[Reversal at any population size]
\label{rem:blocks}
The decay in Table~\ref{tab:reversal} reflects mixing, not size. Because
\eqref{eq:system} is local, a block-diagonal $A$ decouples: each block solves
its own system and $\bth=\big(m^{-1}\sum_ac_a\big)^{-1}$ is the reciprocal of
the population-average capacity. Replicating an identical block leaves that
average, and hence $\bth$, unchanged. Taking $k$ disjoint copies of the network
of Proposition~\ref{prop:reversal} therefore gives $\bth(\lambda)$ independent
of $k$, so that $\bth(1)=0.837209<6/7=\bth(0)$ at $m=3k$ for every $k$; we verify
this numerically up to $m=3000$. The construction is robust to weak ties across
blocks: mixing each row with a fraction $\epsilon$ of the uniform complete
network preserves the reversal at $m=30$ for $\epsilon\leq0.1$
($\bth(1)=0.8526$ at $\epsilon=0.1$) and destroys it by $\epsilon=0.2$
($\bth(1)=0.8681$). Table~\ref{tab:reversal} should therefore be read as
describing a dense ensemble in which local benchmarks approach a common
average, and not as a claim about large networks in general: block structure,
or any persistent correlation between type and network position, sustains the
reversal at arbitrary population size.

Note the contrast with Proposition~\ref{prop:homophily}, which uses the same
block structure to the opposite end. There risk aversion is constant
\emph{within} blocks and the mechanism vanishes; here it is heterogeneous
within blocks and absent \emph{across} them, and the mechanism survives at any
scale. What separates the two is where the heterogeneity sits relative to the
partition, not the topology.
\end{remark}

\subsection{Homophily blocks the mechanism}

\begin{proposition}
\label{prop:homophily}
Partition the investors into groups such that $A_{aj}=0$ whenever $a$ and $j$
belong to different groups, and suppose $\theta_a$ is constant within each
group. Then $c_a=1/\theta_a$ solves \eqref{eq:system} for every profile
$(\lambda_1,\dots,\lambda_m)$, and $\bth$ equals the harmonic mean of $\theta$
regardless of the intensity of comparison.
\end{proposition}

\begin{proof}
Row $a$ of \eqref{eq:system} reads
$(\theta_a+\lambda_a)c_a-\lambda_a\sum_jA_{aj}c_j=1$. If $c_j=1/\theta_j$ and
$\theta_j=\theta_a$ for every $j$ with $A_{aj}>0$, the left-hand side is
$(\theta_a+\lambda_a)/\theta_a-\lambda_a/\theta_a=1$. Uniqueness gives the
claim, and $1/\theta_M=\sum_a1/\theta_a$ is the harmonic mean up to the factor
$m$.
\end{proof}

With two equally sized groups at $\theta\in\{0.5,2\}$, $\bth=0.8$ at
$\lambda=0$, at $\lambda=1$ and at $\lambda=1000$; the arithmetic mean $1.25$
that Proposition~\ref{prop:mf} would predict is never approached. Assortative
matching on risk aversion therefore removes the comparison channel entirely,
which is worth knowing because assortativity on wealth and on financial
sophistication is a documented feature of social networks.

%=====================================================================
\section{An aside on projection beliefs}
\label{sec:projection}

Assumption~\ref{ass:nash} has investors conjecture that others do not respond
to their own trade. A different behavioural assumption, sometimes invoked
under the heading of the false consensus effect \cite{ross}, is that an
investor believes others hold what he holds. This is not a variant of Nash
behaviour --- the conjectured positions are not the equilibrium ones, and are
in general not realised --- and we treat it separately.

Suppose investor $a$ believes that every $j$ in his benchmark holds the same
portfolio \emph{weights} as he does, so that
$\hat{\vec\pi}_{j\mid a}=(W_j^0/W_a^0)\vec\pi_a$. Substituting in
\eqref{eq:pref} and optimising gives
%
\begin{equation}
\kappa_a\vec\pi_a=\Sigma^{-1}\vec r_e,
\qquad
\kappa_a=\theta_a+\lambda_a\left(1-\frac{\bar W_a^0}{W_a^0}\right),
\qquad \bar W_a^0=\sum_jA_{aj}W_j^0 ,
\label{eq:projection}
\end{equation}
%
provided $\kappa_a\neq0$. The coefficient now depends on the investor's
position in the wealth distribution: an investor poorer than his benchmark has
$\kappa_a<\theta_a$ and behaves as though less risk averse.

Two observations. First, the units of the projection decide the outcome: an
investor who projects his \emph{monetary} position rather than his weights
believes $\hat{\vec\pi}_{j\mid a}=\vec\pi_a$, which returns $\kappa_a=\theta_a$
and neutrality. The psychological description ``he assumes others are like
him'' does not by itself determine the equilibrium. Second, $\kappa_a$ may
vanish or turn negative for investors sufficiently poor relative to their
benchmark. The \emph{perceived} individual optimisation problem remains well posed for
$\kappa_a\neq0$ --- the exponent is strictly convex and the optimum interior,
though optimality is relative to a benchmark that is not the one his peers
actually realise --- but the position
is then a short one, and the aggregation \eqref{eq:clear} mixes terms of both
signs, so that $\sum_a\kappa_a^{-1}$ can vanish or change sign and no vector of
expected returns clears the market with a positive premium. This is a
restriction on admissible configurations rather than a failure of individual
optimisation.

We do not develop this case further. It rests on beliefs that are not realised
in equilibrium, and a satisfactory treatment would require a decision-theoretic
statement of what the investor maximises when his perceived benchmark differs
from the positions his peers actually choose, together with a welfare criterion
that does not presume the beliefs are correct.

%=====================================================================
\section{Conclusion}

Where the aggregation conditions of \cite{gomez} hold, relative wealth concerns
in a CARA economy are neutral: with common risk aversion every investor holds
the position he would hold without them, whatever the network of influence and
whatever the dispersion of comparison intensities. Where the conditions fail,
because risk aversion is heterogeneous, the network enters and the effect is
governed by the interaction between preference dispersion and network
structure. Under complete mixing the comparison interpolates the aggregation of
risk aversion from the harmonic mean to the arithmetic mean, at a rate given by
the squared coefficient of variation of shifted risk tolerance. Away from it,
the direction is not guaranteed: a small directed network can make the relation
non-monotone and eventually drive the coefficient below its no-comparison
value. Under a dense i.i.d.-weight design this becomes rare as the population
grows, but that reflects the approach to complete mixing rather than size as
such: with block structure the same reversal survives at any number of
investors. Assortative matching on risk aversion can suppress the effect altogether.

The practical reading is that a representative-agent treatment of relative
concerns is safe exactly where Gómez's conditions hold and unreliable
elsewhere, in a way that depends on features of the interaction structure that
a representative agent cannot express. Characterising the classes of networks
for which the mean-field direction is guaranteed --- reversibility, double
stochasticity, bounds on assortativity between risk aversion and centrality ---
would make the result usable, and is the obvious next step.

%=====================================================================
\section*{Replication}
\label{sec:replication}

All numerical results are reproducible from the repository
%
\begin{center}
\url{https://github.com/renatovicente/price-of-risk-aggregation}
\end{center}
%
The module \texttt{scripts/network.py} solves \eqref{eq:system};
\texttt{scripts/neutrality.py} verifies Theorem~\ref{thm:neutral} and
Proposition~\ref{prop:mf}; \texttt{scripts/sign\_reversal.py} produces
Table~\ref{tab:reversal}, including the network and type generators and the
seeds; \texttt{scripts/block\_reversal.py} produces the block replication and
the mixing robustness of Remark~\ref{rem:blocks}. The scripts depend only on
\texttt{numpy} and \texttt{scipy}, read no data and draw from fixed seeds;
\texttt{./reproduce.sh} runs all four in about fifteen seconds.

%=====================================================================
\begin{thebibliography}{99}

\bibitem{abel} A. B. Abel,
\emph{Asset prices under habit formation and catching up with the Joneses},
American Economic Review Papers and Proceedings 80 (1990) 38--42.

\bibitem{ballester} C. Ballester, A. Calv\'o-Armengol, Y. Zenou,
\emph{Who's who in networks. Wanted: the key player},
Econometrica 74 (2006) 1403--1417.

\bibitem{campbell} J. Y. Campbell, J. H. Cochrane,
\emph{By force of habit: a consumption-based explanation of aggregate stock
market behavior}, Journal of Political Economy 107 (1999) 205--251.

\bibitem{chankogan} Y. L. Chan, L. Kogan,
\emph{Catching up with the Joneses: heterogeneous preferences and the dynamics
of asset prices}, Journal of Political Economy 110 (2002) 1255--1285.

\bibitem{degennaro} L. De Gennaro Aquino, E. G. De Giorgi, Y. Lou, M. S. Strub,
\emph{Investment and asset pricing with relative wealth concerns and multiple
risky assets},
The North American Journal of Economics and Finance 84 (2026) 102628.

\bibitem{demarzo04} P. DeMarzo, R. Kaniel, I. Kremer,
\emph{Diversification as a public good: community effects in portfolio choice},
Journal of Finance 59 (2004) 1677--1716.

\bibitem{dupor} B. Dupor, W.-F. Liu,
\emph{Jealousy and equilibrium overconsumption},
American Economic Review 93 (2003) 423--428.

\bibitem{gali} J. Gal\'i,
\emph{Keeping up with the Joneses: consumption externalities, portfolio choice,
and asset prices},
Journal of Money, Credit and Banking 26 (1994) 1--8.

\bibitem{gollier} C. Gollier,
\emph{Wealth inequality and asset pricing},
Review of Economic Studies 68 (2001) 181--203.

\bibitem{gomez} J.-P. G\'omez,
\emph{The impact of keeping up with the Joneses behavior on asset prices and
portfolio choice},
Finance Research Letters 4 (2007) 95--103.

\bibitem{kraft} H. Kraft, A. Meyer-Wehmann, F. T. Seifried,
\emph{Dynamic asset allocation with relative wealth concerns in incomplete
markets},
Journal of Economic Dynamics and Control 113 (2020) 103857.

\bibitem{lacker} D. Lacker, T. Zariphopoulou,
\emph{Mean field and $n$-agent games for optimal investment under relative
performance criteria},
Mathematical Finance 29 (2019) 1003--1038.

\bibitem{panageas} S. Panageas,
\emph{The implications of heterogeneity and inequality for asset pricing},
NBER Working Paper 26974 (2020).

\bibitem{ross} L. Ross, D. Greene, P. House,
\emph{The false consensus effect: an egocentric bias in social perception and
attribution processes},
Journal of Experimental Social Psychology 13 (1977) 279--301.

\end{thebibliography}

\end{document}
